\chapter{Context-Sensitive Analysis}

\section{Beyond context-free}

Many rules of a programming language cannot be expressed by a CFG (or would make it
hopelessly complicated):
\begin{itemize}
  \item a variable must be declared before use, and not declared twice in the same scope;
  \item a procedure must be called with the right number and types of arguments;
  \item operands must have compatible types.
\end{itemize}
The language $\{ wcw : w\in\{a,b\}^*\}$ (``declare $w$, then use $w$'') is not context-free.
So after parsing, the compiler performs \textbf{context-sensitive (semantic) analysis} by
traversing the parse tree with auxiliary data structures. A program that passes all checks is
a valid WLP4 program; only valid programs are given to the code generator.

\section{The WLP4 language}

WLP4 is a tiny subset of C++:
\begin{itemize}
  \item Two types: \code{int} (32-bit signed) and \code{int*} (pointer to int).
  \item Procedures return \code{int}; the last procedure is \code{wain}, which takes two
        parameters (the second must be \code{int}) and is the entry point.
  \item All local declarations appear at the top of a procedure, before any statements, and
        must be initialised: \code{int x = 5;} or \code{int* p = NULL;}.
  \item Statements: assignment; \code{if} with a mandatory \code{else} (braces required
        around both branches); \code{while}; printing with \code{println}; and freeing
        with \code{delete []}.
  \item Expressions: \code{+ - * / \%}, comparisons \code{== != < <= > >=} (only in tests),
        address-of \code{\&}, dereference \code{*}, \code{new int[n]}, \code{NULL}, and
        procedure calls.
  \item No globals, no nested procedures; a procedure must be declared before it is called.
\end{itemize}

\begin{lstlisting}[style=cpp]
int sum(int* a, int n) {
    int i = 0;
    int s = 0;
    while (i < n) {
        s = s + *(a + i);
        i = i + 1;
    }
    return s;
}
int wain(int* a, int n) {
    return sum(a, n);
}
\end{lstlisting}

\section{Symbol tables}

\begin{definition}
A \textbf{symbol table} maps identifiers to their attributes: for variables, the type (and
later the stack offset); for procedures, the signature (parameter types).
\end{definition}

WLP4 has two levels of scope:
\begin{itemize}
  \item a \textbf{global} table of procedures: name $\mapsto$ (signature, local table);
  \item one \textbf{local} table per procedure: parameter and variable names $\mapsto$ type.
\end{itemize}
A variable and a procedure may share a name; inside a procedure a local variable shadows a
procedure name for the purpose of calls (\code{f()} where \code{f} is a local variable is an
error).

\paragraph{Building it.} Traverse the parse tree:
\begin{itemize}
  \item At \code{procedure}: error if the name is already in the global table; create a new
        local table; collect the parameter types into the signature.
  \item At each \code{dcl}: error if the name is already in the current local table;
        otherwise insert it with its type (\code{type -> INT} gives \code{int},
        \code{type -> INT STAR} gives \code{int*}).
  \item At each use (\code{factor -> ID}, \code{lvalue -> ID}): error if not in the current
        local table.
  \item At each call \code{ID(...)}: error if the name is a local variable or not a
        previously declared procedure (recursion is allowed, so the current procedure is in
        the table as soon as its signature is read).
\end{itemize}

\begin{lstlisting}[style=cpp]
struct ProcInfo { std::vector<Type> params; std::map<std::string, Type> locals; };
std::map<std::string, ProcInfo> procs;

void collectDcl(Tree *dcl, ProcInfo &p) {
    std::string name = dcl->children[1]->lexeme;
    Type t = (dcl->children[0]->children.size() == 2) ? INTSTAR : INT;
    if (p.locals.count(name)) throw SemanticError("duplicate " + name);
    p.locals[name] = t;
}
\end{lstlisting}

\section{Type checking}

We compute a type for each expression bottom-up (a post-order traversal), and check that
statements are \textbf{well-typed}. The rules can be written as inference rules; here as a
table.

\subsection{Expressions}
\begin{center}
\small
\begin{tabular}{lll}
\toprule
expression & requirement & type \\
\midrule
\code{NUM} & & \code{int} \\
\code{NULL} & & \code{int*} \\
\code{ID} & declared & declared type \\
\code{(E)} & & type of $E$ \\
\code{\&E} & $E$ an lvalue of type \code{int} & \code{int*} \\
\code{*E} & $E : $ \code{int*} & \code{int} \\
\code{new int[E]} & $E :$ \code{int} & \code{int*} \\
\code{E1 * E2}, \code{/}, \code{\%} & both \code{int} & \code{int} \\
\code{E1 + E2} & \code{int + int} & \code{int} \\
 & \code{int* + int} or \code{int + int*} & \code{int*} \\
\code{E1 - E2} & \code{int - int} & \code{int} \\
 & \code{int* - int} & \code{int*} \\
 & \code{int* - int*} & \code{int} \\
\code{f(E1,...,En)} & arity and each type match the signature & \code{int} \\
\bottomrule
\end{tabular}
\end{center}

The pointer rules mirror C: adding an integer to a pointer moves it by that many
\emph{elements}; subtracting two pointers gives the distance in elements. \code{int* + int*}
is an error.

\subsection{Statements and tests}
\begin{itemize}
  \item \code{lvalue = expr;} : both sides have the same type.
  \item \code{println(E);} : $E$ is \code{int}.
  \item \code{delete [] E;} : $E$ is \code{int*}.
  \item tests \code{E1 op E2} : $E_1$ and $E_2$ have the same type (comparing pointers is
        allowed, comparing a pointer to an int is not).
  \item \code{if}/\code{while}: the test is well-typed and the bodies are well-typed.
\end{itemize}

\subsection{Declarations and procedures}
\begin{itemize}
  \item \code{int x = NUM;} requires \code{int}; \code{int* p = NULL;} requires \code{int*}.
  \item The \code{return} expression of every procedure (including \code{wain}) must be \code{int}.
  \item The second parameter of \code{wain} must be \code{int}.
\end{itemize}

\begin{lstlisting}[style=cpp]
Type typeOf(Tree *t) {
    const std::string &r = t->rule;
    if (r == "factor NUM")  return INT;
    if (r == "factor NULL") return INTSTAR;
    if (r == "factor ID")   return lookup(t->children[0]->lexeme);
    if (r == "factor STAR factor") {
        if (typeOf(t->children[1]) != INTSTAR) throw SemanticError("deref non-pointer");
        return INT;
    }
    if (r == "expr expr PLUS term") {
        Type a = typeOf(t->children[0]), b = typeOf(t->children[2]);
        if (a == INT && b == INT) return INT;
        if (a != b) return INTSTAR;        // exactly one pointer
        throw SemanticError("int* + int*");
    }
    // ... remaining rules ...
}
\end{lstlisting}

It is convenient to store the computed type in each tree node, because the code generator
needs it later (e.g.\ to scale pointer arithmetic by 4).

\section{Attribute grammars (brief)}
\begin{aside}
More generally, semantic analysis can be described by an \emph{attribute grammar}: each
grammar symbol carries attributes (like \emph{type}) and each production has rules for computing
them. \emph{Synthesized} attributes flow up the tree (types of expressions); \emph{inherited}
attributes flow down (the current symbol table). WLP4 type checking uses mostly synthesized
attributes plus a symbol table passed down.
\end{aside}

\section{Errors and pitfalls}
\begin{itemize}
  \item Check both ``declared twice'' and ``used but not declared'', and do it per procedure.
  \item Recursion: insert the procedure into the global table before checking its body.
  \item A call to an undeclared (or later-declared) procedure is an error: WLP4 requires
        declaration before use.
  \item Remember \code{\&} needs an lvalue: \code{\&(x)} is fine (parenthesised lvalue), \code{\&5} is not
        parseable anyway.
  \item Do not stop at the first type of each subtree: the whole tree must be checked, even
        statements whose value is never used.
\end{itemize}

\section{Summary}
\begin{itemize}
  \item Context-sensitive rules (declarations, types, arity) are checked by traversing the parse tree.
  \item Symbol tables: one global table of procedure signatures and one local table per procedure.
  \item WLP4 types are \code{int} and \code{int*}; pointer arithmetic follows C rules.
  \item Compute expression types bottom-up; check statements, tests, declarations and returns.
  \item Annotate the tree with types; the code generator will use them.
\end{itemize}
